\begin{statement}{E.30}{Proved}{A necessary finite conjugacy test from the four-strand periodic involution}
Let $q:B_4\twoheadrightarrow T$ be any finite quotient of the
discrete four-strand braid group, and write
\[
 a=q(\sigma_1),\quad b=q(\sigma_2),\quad c=q(\sigma_3),
 \quad d=abc,\quad z=d^4,\quad C=\langle z\rangle.
 \tag{E.30.1}
\]
The subgroup $C$ is central.  For every
$F=(\lambda,f)\in\widehat{\mathrm{GT}}$, let
\[
 a_F=q(F(\sigma_1)),\qquad
 b_F=q(F(\sigma_2)),\qquad
 c_F=q(F(\sigma_3)),
 \tag{E.30.2}
\]
where $q$ denotes also its continuous extension and $F$ acts by
the usual based Drinfeld braid formulas.  Then
\[
       (a_Fb_Fc_F)^2C\sim_{T/C}d^2C.
 \tag{E.30.3}
\]
Equivalently, the finite equation
\[
       (a_Fb_Fc_F)^2=h d^2h^{-1}z^j
       \quad\text{has a solution }h\in T,\quad
                       j\in\mathbb Z/|z|.
 \tag{E.30.4}
\]
This condition applies to the actual Ferrari quotient maps used
in Part E.  It retains the complete conjugacy class in $T/C$;
the full-twist equation alone prescribes only the square of the
element on the left of \textup{(E.30.3)}.  The assertion uses
the homomorphism $q\circ F:\widehat B_4\to T$ and requires no
invariance of $\ker q$ under $F$.
\end{statement}
\begin{proof}
We give the group comparison that brings the torsion theorem
into this particular finite target.  Capping the boundary of
a four-punctured disc by a once-punctured disc gives the exact
discrete homomorphism
\[
 B_4/\langle(\sigma_1\sigma_2\sigma_3)^4\rangle
       \xrightarrow{\ \simeq\ }
 \Gamma_{0,[5]}^{(5)}\ \leq\ \Gamma_{0,[5]},
 \tag{E.30.5}
\]
where the subgroup on the right fixes the fifth marked point
and has index five.  Indeed a mapping class fixing that point
can be represented by a map fixing a small disc about it;
restriction to the complementary four-punctured disc gives a
braid.  Changing the isotopy at the capped disc changes this
braid by an integral boundary twist, precisely the central
full twist displayed on the left.  This identifies both the
kernel and the image in \textup{(E.30.5)}.

Completion of the quotient on the left is
$\widehat B_4/\overline{\langle(\sigma_1\sigma_2\sigma_3)^4\rangle}$:
their finite quotients are exactly the finite braid quotients
killing that central element.  A finite-index subgroup has
the full induced profinite topology, by taking finite-index
cores.  Thus completion of \textup{(E.30.5)} identifies its
left side with the open fifth-point stabilizer in
$\widehat\Gamma_{0,[5]}$.
For the action comparison, use the specified homomorphisms
\[
 \widehat B_4\xrightarrow{\ \iota\ }\widehat B_5
       \xrightarrow{\ \pi\ }\widehat\Gamma_{0,[5]},
 \qquad \iota(\sigma_i)=\sigma_i\quad(1\leq i\leq3).
 \tag{E.30.5a}
\]
Here $\pi$ kills the closed normal subgroup generated by
\[
 z_5=(\sigma_1\sigma_2\sigma_3\sigma_4)^5,
 \qquad y_5=\sigma_4\sigma_3\sigma_2\sigma_1^2
                         \sigma_2\sigma_3\sigma_4.
 \tag{E.30.5b}
\]
These are exactly the central full twist and the sphere
Hurwitz relation in Collas's presentation
\cite[Section 3.1 and Lemma 3.12]{ERootCollas2012}.
The composite $\pi\iota$ is the capping map in
\textup{(E.30.5)}: it places the fifth puncture outside the
four-strand disc and then fills its outer boundary.
The based Drinfeld formula for each of the first three
generators uses the same generators and the same preceding
pure braid in $B_4$ and $B_5$, so $F\iota=\iota F$ on
those generators and hence on the completed group.
Collas's Lemma 3.12 gives the literal equalities
$F(y_5)=y_5^\lambda$ and $F(z_5)=z_5^\lambda$.
They make the displayed normal subgroup invariant and give
$\pi F=F\pi$.  Thus \textup{(E.30.5a)} proves equivariance
of the actual capping map and its fifth-point stabilizer.

Let $\delta=\sigma_1\sigma_2\sigma_3$ and let $\epsilon$ be
the image of $\delta^2$ in \textup{(E.30.5)}.  It has order
exactly two: its square is the killed full twist, and its
permutation on five marked points is a product of two disjoint
transpositions with the unique fixed point $5$.
Collas's Theorem~3.10, also stated as Theorem~B in the
introduction, gives a conjugator
$u\in\widehat\Gamma_{0,[5]}$ with
\[
                         F(\epsilon)=u\epsilon^\lambda u^{-1}
                                      =u\epsilon u^{-1};
 \tag{E.30.6}
\]
the second equality uses $\lambda\in\widehat{\mathbb Z}^{\times}$
and the order two of $\epsilon$
\cite{ERootCollas2012}.
The theorem concerns the full profinite GT group, with its
genus-zero braid action, so its source is exactly the source
in the present statement.

We now locate the conjugator.  The Drinfeld action preserves
the permutation of each braid generator: its exponent is a
unit, hence odd, and its conjugating word belongs to the
commutator subgroup of a pure free subgroup.  Consequently
$F(\epsilon)$ has the same permutation as $\epsilon$.
The permutation of $u$ in \textup{(E.30.6)} sends the unique
fixed point of the latter permutation to the unique fixed
point of the former.  It fixes $5$.  Hence $u$ lies in the
open subgroup in \textup{(E.30.5)} itself.  We have obtained
the required conjugacy inside the completed four-strand
central quotient.

The homomorphism $\widehat B_4\to T/C$ kills the closed
central subgroup in \textup{(E.30.5)}, so it factors through
that quotient.  Applying it to \textup{(E.30.6)} gives
\textup{(E.30.3)}.  Lifting the conjugator to $T$ gives
\textup{(E.30.4)}, since the two sides then differ by an
element of $C$.  Every map in this argument is either the
specified quotient map, a finite-index inclusion and its
completion, or the stated GT action.  In particular no
endomorphism of the particular finite quotient $T$ has been
assumed.  For a finite candidate, failure of
\textup{(E.30.4)} excludes that candidate from this full-GT
image; a conclusion about a whole CM fibre still requires
the complete fibre and multiplier comparisons.
\end{proof}
